\subsection{Feature Ablation Study}
\label{sec:feature_ablation}
\label{sec:feature_ablation_study}

The quantification of classification performance from single-domain toward dual- and tri-domain configurations is summarized in Table~\ref{tab:VII}, Table~\ref{tab:VIII}, Table~\ref{tab:IX}, and Table~\ref{tab:X}. In the SVM model, representation transition from single-domain Face (90.83\%) to dual-domain Emotion $\oplus$ Face (93.29\%) increases accuracy by 2.46\%. Adding a third domain in the tri-domain Face $\oplus$ Emotion $\oplus$ Age configuration further elevates accuracy to 93.70\% (cumulative +2.87\%), with its 95\% CI [92.60\%, 94.65\%] overlapping the dual-domain interval [92.15\%, 94.27\%]. The close consistency of the CV Score (92.65\% $\pm$ 0.93\%) with test accuracy confirms decision boundary stability without overfitting. This progressive improvement demonstrates that combining complementary visual representations reliably sharpens intersectional demographic separation.

An analysis of the age domain (Age) examines the role of aging signals within the demographic classification system. Across all classifiers, the Age feature yields the lowest accuracy as a single domain, recording 87.64\% for SVM and 69.63\% for GNB. The performance decline in Random Forest (-0.65\% from dual-domain to tri-domain) occurs because expanded feature dimensions make random trees prone to choosing less informative age features at split points. Conversely, the SVM model with a polynomial kernel simultaneously evaluates interactive feature dimensions. This finding refutes the premise that adding features unconditionally improves accuracy; rather, it depends on classifier characteristics in handling feature spaces, confirming that feature representations reach an empirical saturation point.
